\section{Architecture of Nested Persistence}
\label{sec:nested_persistence}

A persistent system must service its continuing demands within finite resources while retaining the distinctions needed for its future operation. When a single control loop cannot meet a fixed workload within its causal deadline, the architecture must change: it may acquire additional parallel capacity, reduce the workload, or relax the deadline. Partitioning is required when the workload and deadline are fixed and additional capacity is supplied by parallel subsystems. Dividing an unchanged, insufficient budget cannot itself remove the overload.

We formulate three architectural rules: Capacity Partitioning accounts for competing demands and interface costs; Reversible Encoding preserves distinctions lost by a coarse description; and Impedance Matching controls the resource burden of transfer between subsystems. Their mathematical statements below are independent of the vacuum example. They constrain which architectures are admissible. A unique allocation additionally requires the workload, available operations, cost law, and selection objective of the application.

For example, a firm allocating a fixed budget among manufacturing, innovation, and compliance must include coordination costs and a reserve. An indivisible remainder may be retained in its accounts rather than erased. A communication interface must deliver the required throughput within its processing budget. These examples motivate the rules; they do not assume that an organizational interface obeys the particular wave-impedance law used later for the vacuum.

\subsection{Capacity Partitioning}
\label{rule_capacity_partitioning}

\textbf{Persistence constraint.} Finite regulatory variety and finite execution rates motivate a distinction between the number of states a controller can represent and the work it can perform before a deadline \cite{raadt_ashbys_1987,margolus_maximum_1998}. Neither quantity should be substituted for the other without a model connecting them.

Let $w_i\geq0$ be the work demanded by task $i$ during a specified interval, $C_j$ the work capacity of subsystem $j$ in that interval, and $h_j\geq0$ its interface overhead. A feasible allocation assigns nonnegative service $s_{ij}$ satisfying
\begin{equation}
    \sum_j s_{ij}=w_i,
    \qquad
    \sum_i s_{ij}+h_j\leq C_j.
\end{equation}
Summing these constraints gives the necessary global budget condition
\begin{equation}
    \sum_i w_i+\sum_j h_j\leq\sum_j C_j.
\end{equation}
For equal subsystem capacity $C>0$, this implies that the number of subsystems is at least $(\sum_iw_i+\sum_jh_j)/C$. This aggregate bound is necessary; dependencies, indivisible tasks, or communication deadlines may impose further scheduling restrictions.

\textbf{Operational mandate.} The system must keep its demand and coordination cost within available capacity. It need not consume every available unit. Writing used capacity plus reserve equal to total capacity is an accounting identity; setting the reserve to zero is an additional saturation condition. Distinct allocations can obey the same feasibility constraints. When a unique split is claimed, its selection objective must be stated and proved to select that split.

\textbf{Formal scope.} Lean proves the total budget inequality, the impossibility of allocating an insufficient unchanged budget, and the corresponding worker-count bound. It also exhibits a family of feasible allocations and a feasible allocation with idle capacity (\texttt{General.Capacity}). The model treats work as divisible and assumes that all specified demand must be completed within the interval.

\paragraph{Repeated workloads and bounded backlog.}
Work may vary within a repeating cycle. Write $d_{i,t}$ and $s_{i,t}$ for demand
and service of class $i$ at phase $t$ in a cycle of $N$ phases, with common
operation rate $\rho$. A queue obeying
\begin{equation}
 q_{i,n,t+1}=q_{i,n,t}+\rho d_{i,t}-s_{i,t},\qquad
 q_{i,n+1,0}=q_{i,n,N}
\end{equation}
accumulates the same net drift in each cycle. If its nonnegative boundary backlog
is bounded for all cycles, this drift must vanish:
\begin{equation}
 \sum_t s_{i,t}=\rho\sum_t d_{i,t}.
\end{equation}
For $\rho\ne0$, normalized cycle-total service shares therefore equal normalized
cycle-total demand shares (Lean: \texttt{General.Service}). This permits deferred
service within the cycle; it does not require pointwise equality of demand and
service, uniform phases, or full resource utilization. The repeating demand and
service pattern, the queue balance law, and bounded backlog are hypotheses.
This necessary balance does not itself construct a schedule meeting every
causal deadline.

\subsection{Reversible Encoding}
\label{rule_reversible_encoding}

\textbf{Persistence constraint.} A coarse description may merge states that remain physically distinguishable. If those distinctions must be conserved within the system, the information omitted by the description must remain recoverable elsewhere. Reversible computation provides an operational formulation of this requirement \cite{bennett_thermodynamics_1982}.

For a state space $X$, let $q:X\to Q$ be the coarse description and $m:X\to M$ its retained metadata. The combined encoding is
\begin{equation}
    E(x)=(q(x),m(x)).
\end{equation}
Conservation of the specified distinctions requires $E$ to be injective. Equivalently, the metadata must distinguish every pair of states merged by $q$; equivalently, there is a decoder that recovers every encoded state. For finite state spaces,
\begin{equation}
    |q^{-1}(a)|\leq |M| \quad\text{for every coarse label }a,
    \qquad |X|\leq |Q|\,|M|.
\end{equation}
Thus, if a quotient label admits $k$ distinct possible remainders, an encoding retaining all of them requires at least $k$ metadata states.

\textbf{Operational mandate.} Preserve the distinctions required by the dynamics, either in the primary description or in recoverable metadata. This is an information-capacity condition. Identifying the metadata with physical mass, or assigning a thermodynamic cost to erasure, requires the application's additional physical model. Landauer's principle supplies thermodynamic context \cite{landauer_irreversibility_1961}; the counting theorem does not itself assert an infinite heat cost for a finite erasure.

\textbf{Formal scope.} Lean proves the reversible-encoding criterion, its equivalent decoder formulation, the fiber and total capacity bounds, and invariance under relabeling the metadata (\texttt{General.Encoding}).

For group-valued data on sites, an exact relative description consists of the
reference value $x_o$ together with $x_i-x_o$ for $i\ne o$. Lean supplies an
explicit inverse (\texttt{General.RelationalEncoding}). The self-difference
vanishes, but dropping the reference value merges distinct constant
configurations. A reference position alone does not supply that value. This
information-preservation statement does not set the maintenance frequency or
cost of the retained data.

\subsection{Impedance Matching}
\label{rule_impedance_matching}

\textbf{Persistence constraint.} An interface must deliver the required transfer while keeping its processing and reflected-load burden within the resources available to handle it. Exact matching is a possible optimum, but a nonzero mismatch does not by itself establish failure of persistence. The transfer law and the available reserve must be specified.

For the positive-real interface law considered here, source and load impedances $Z_s,Z_l>0$ give
\begin{equation}
    r=\frac{Z_l-Z_s}{Z_l+Z_s},
    \qquad T=1-r^2.
\end{equation}
An incident load $P>0$ divides into reflected load $Pr^2$ and transmitted load $PT$, with
\begin{equation}
    Pr^2+PT=P.
\end{equation}
Reflection remains in this conserved ledger; it is not identified with irreversible erasure. If the model allocates reserve $B$ to handling reflected load, admissibility requires
\begin{equation}
    Pr^2\leq B
    \quad\Longleftrightarrow\quad
    T\geq 1-\frac{B}{P}.
\end{equation}
For zero reserve, exact matching $Z_l=Z_s$ is necessary. With positive reserve, mismatched interfaces can remain admissible. Exact matching uniquely maximizes transmission for this law, but that optimality statement is stronger than budget feasibility.

\textbf{Operational mandate.} Derive the tolerated mismatch from the interface's resource budget and transfer law. A tolerance may not be chosen merely to cover a numerical residual. Conversely, a numerical near match must not be described as an exact identity. If a physical application uses a different transfer or recovery mechanism, that mechanism must supply the corresponding admissibility condition.

\textbf{Formal scope.} Lean proves the conserved ledger, the reflected-load budget criterion, exact matching under zero reserve, and an explicit mismatched interface within a positive reserve (\texttt{General.Matching}). This is a conditional interface model, not a claim that every persistent system obeys the displayed impedance formula.

\subsection{Recursive Universality}
\label{sec:recursive_universality}

The rules apply recursively when each subsystem must service demands, preserve distinctions, and exchange information across a boundary. They describe the resource conditions under which nearly decomposable architectures can persist \cite{klir_architecture_1991}. Their application to a particular substrate must identify the resources and operations at each level.

Finite total budget alone does not prove finite nesting depth if successive levels may cost arbitrarily little. If every level costs at least $c_{\min}>0$ and a branch has budget $B$, then
\begin{equation}
    n\,c_{\min}\leq B,
    \qquad n\leq\frac{B}{c_{\min}}.
\end{equation}
Alternatively, a natural-valued resource that strictly decreases at every level bounds the depth by its initial value. Lean proves both bounds and gives a counterexample to omitting the uniform cost floor: the positive costs $2^{-(i+1)}$ have every finite prefix sum within budget one (\texttt{General.Hierarchy}). Finite branching is additionally needed to turn a depth bound into a finite-total-node statement.

The constants example that follows must therefore supply its own workload, encoding, response law, and selection criterion. The general rules constrain those choices; they do not, on their own, determine unique numerical partition coefficients. The Lean general-rules layer imports no $E_8$ or vacuum-constant module, and a separate dependency audit checks that separation.
